% =============================================================================
%  Glossaire
%
%  Destiné au lecteur non spécialiste. Il ne remplace pas les encadrés
%  « Définition » des chapitres — qui définissent chaque notion à son premier
%  emploi, dans son contexte — mais permet de retrouver une définition sans
%  relire le chapitre où elle apparaît.
%
%  À placer de préférence APRÈS la liste des acronymes, en fin de pages
%  liminaires (ou en annexe si le modèle le prévoit ainsi).
% =============================================================================

\chapter*{Glossaire}
\addcontentsline{toc}{chapter}{Glossaire}
\markboth{GLOSSAIRE}{GLOSSAIRE}

\begin{description}\setlength{\itemsep}{0.35em}

\item[Actif] Objet dans lequel on peut placer de l'argent~: action d'entreprise,
obligation d'État, fonds indiciel.

\item[Allocation] Répartition d'un capital entre plusieurs actifs. Synonyme
usuel~: \emph{poids} du portefeuille.

\item[Backtest] Simulation d'une stratégie sur des données passées, en rejouant
le temps tel qu'il s'est écoulé.

\item[Biais de survivance / de vision du futur] Toute situation où une décision
passée est influencée par une information qui n'était pas disponible à ce
moment-là. C'est l'erreur la plus grave possible dans ce domaine~: elle produit
des performances impossibles à reproduire en conditions réelles.

\item[Bootstrap à blocs] Méthode d'estimation de l'incertitude consistant à
reconstituer des milliers d'historiques fictifs en retirant au sort des blocs de
journées consécutives, afin de préserver la structure temporelle des rendements.

\item[Coefficient d'information] Corrélation de rang entre rendements prédits et
rendements réalisés~: mesure la capacité d'un signal à \emph{classer} les actifs,
sans exiger que ses prédictions soient justes en niveau.

\item[Data snooping] Biais produit lorsque de nombreuses stratégies sont
essayées sur les mêmes données et que seule la meilleure est retenue. Même sans
fraude ni fuite temporelle, le maximum observé est alors optimiste.

\item[Corrélation] Mesure du degré auquel deux actifs évoluent ensemble. Comprise
entre $-1$ (mouvements opposés) et $+1$ (mouvements identiques).

\item[Covariance (matrice de)] Tableau décrivant conjointement la volatilité de
chaque actif et la corrélation de chaque paire. C'est la description du risque
d'un portefeuille, et la quantité la plus difficile à estimer.

\item[\emph{Drawdown} (perte maximale)] Baisse maximale subie entre un sommet et
le creux qui le suit. Répond à la question « combien aurais-je pu perdre au pire
moment~? ».

\item[Équipondération (1/N)] Stratégie consistant à placer la même somme sur
chaque actif. Sans paramètre à estimer, donc sans erreur d'estimation~: c'est la
référence de comparaison honnête.

\item[Fonds indiciel (ETF)] Fonds coté en bourse répliquant un indice de marché~;
en acheter une part revient à détenir une fraction de tout l'indice.

\item[Hors échantillon] Se dit d'une performance mesurée sur des données qui
n'ont servi ni à entraîner ni à sélectionner le modèle évalué.

\item[Modèle de Markov caché] Modèle statistique supposant que les observations
sont engendrées par un état non observable qui change au cours du temps~; l'état
est inféré à partir de ce qu'il produit.

\item[Numéraire] Monnaie dans laquelle un patrimoine, son profit et son risque
sont mesurés. Pour un investisseur marocain non couvert, un actif en dollars doit
inclure la variation USD/MAD dans son rendement en dirhams.

\item[Population Stability Index (PSI)] Indicateur comparant la distribution
d'une variable ou d'un score entre une période de référence et une période
évaluée. Il signale une dérive ; il n'en explique pas la cause.

\item[Point de base] Un centième de point de pourcentage ($0{,}01~\%$). Unité
usuelle des coûts de transaction.

\item[Ratio de Calmar] Rendement annualisé divisé par la perte maximale. Mesure
le rendement obtenu par unité de \emph{pire} perte, là où le ratio de Sharpe
divise par la volatilité moyenne.

\item[Ratio de Sharpe] Rendement obtenu par unité de risque, le risque étant
mesuré par la volatilité. Principal indicateur de comparaison entre stratégies —
et, comme le montre le chapitre~\ref{chap:resultats}, un indicateur dont le
choix n'est pas neutre.

\item[Régime de marché] Période durant laquelle le marché se comporte de manière
homogène (par exemple~: hausse régulière et faible volatilité, ou baisse et forte
volatilité). Les régimes ne sont pas observables directement.

\item[Rendement logarithmique] Rendement calculé comme $\ln(P_t/P_{t-1})$. Se
somme dans le temps, contrairement au rendement en pourcentage — d'où son usage
systématique dans ce projet.

\item[Rétrécissement] Aussi appelé \emph{shrinkage}, technique consistant à rapprocher une
estimation bruitée d'une cible simple et stable. On accepte un biais volontaire
en échange d'une forte réduction de variance.

\item[Sharpe déflaté] Probabilité qu'une performance mesurée soit réellement
positive après correction du nombre de stratégies essayées et de la
non-normalité des rendements.

\item[Turnover (rotation)] Proportion du portefeuille effectivement échangée lors
d'un rééquilibrage. Détermine directement les coûts de transaction supportés.

\item[Sélection strictement antérieure] Découpage de validation dans lequel
l'entraînement précède toujours la fenêtre de validation, séparé d'elle par une
purge et un embargo. À la différence de la validation croisée purgée en blocs,
aucune date postérieure à la validation ne peut entrer dans l'entraînement~:
c'est la seule forme qui simule une décision d'allocation séquentielle.

\item[Test pairé] Comparaison de deux stratégies sur les \emph{mêmes} dates, en
rééchantillonnant les deux séries avec les mêmes blocs. Teste la
\emph{différence}, ce que deux intervalles de confiance marginaux ne font pas~:
ils peuvent se chevaucher alors que la différence est systématiquement positive.

\item[Validation croisée purgée] Variante de la validation croisée adaptée aux
séries temporelles~: les observations d'entraînement dont l'horizon chevauche le
bloc de test sont retirées (\emph{purge}), et une bande de journées suivant ce
bloc est écartée (\emph{embargo}).

\item[Volatilité] Ampleur des variations d'un actif, mesurée par l'écart-type de
ses rendements. Mesure usuelle — et imparfaite — du risque.

\end{description}
